\documentclass[11pt,a4paper]{article}
\usepackage[margin=22mm]{geometry}
\usepackage{amsmath,graphicx,booktabs}
\setlength{\parindent}{0pt}
\setlength{\parskip}{6pt}
\newcommand{\evidence}[3]{%
  \begin{center}
    \includegraphics[width=\linewidth,height=#2\textheight,keepaspectratio]{figures/#1.png}\\[4pt]
    {\small #3}
  \end{center}
}
\newcommand{\comparison}[4]{%
  \begin{center}
  \begin{minipage}[t]{0.48\linewidth}\centering
    \includegraphics[width=\linewidth,height=0.48\textheight,keepaspectratio]{figures/#1.png}\\[6pt]
    {\small #2}
  \end{minipage}\hfill
  \begin{minipage}[t]{0.48\linewidth}\centering
    \includegraphics[width=\linewidth,height=0.48\textheight,keepaspectratio]{figures/#3.png}\\[6pt]
    {\small #4}
  \end{minipage}
  \end{center}
}
\begin{document}
\begin{center}
{\Large\bfseries Logic Synthesis and Verification}\\[4pt]
{\large Programming Assignment 1}\\[4pt]
Yung-Hsiang Yang \quad B11901064 \quad September 21, 2026
\end{center}
\section{Exercise 2: two-bit unsigned multiplier}

The file \texttt{mul.blif} implements \(Y=A B\), with
\(A=(a_1a_0)\), \(B=(b_1b_0)\), and \(Y=(y_3y_2y_1y_0)\).
Four partial products and two half-adder stages implement
\[
\begin{aligned}
p_{00}&=a_0b_0,& p_{01}&=a_0b_1,& p_{10}&=a_1b_0,& p_{11}&=a_1b_1,\\
y_0&=p_{00},& y_1&=p_{01}\oplus p_{10},&
c_1&=p_{01}p_{10},\\
y_2&=p_{11}\oplus c_1,& y_3&=p_{11}c_1.
\end{aligned}
\]
Each AND and XOR is expressed as a BLIF \texttt{.names} table.
The input order is \texttt{a1 a0 b1 b0}; the output order is
\texttt{y3 y2 y1 y0}.

An independent reference BLIF was generated from integer multiplication,
rather than the half-adder equations. ABC's \texttt{cec} reported
\texttt{Networks are equivalent}. The reference covers all sixteen input pairs:
\begin{center}
\begin{tabular}{c|rrrr}
\toprule
\(A\backslash B\) & 0 & 1 & 2 & 3\\
\midrule
0 & 0 & 0 & 0 & 0\\
1 & 0 & 1 & 2 & 3\\
2 & 0 & 2 & 4 & 6\\
3 & 0 & 3 & 6 & 9\\
\bottomrule
\end{tabular}
\end{center}
In particular, \(3\cdot3=9=(1001)_2\).

\subsection*{Experimental setup and measured statistics}
The experiments use the course ABC checkout, built with \texttt{make -j4}
and bundled CUDD on an Apple Silicon Mac. Each comparison starts by reading
the same multiplier. Screenshots show actual ABC command output and
ABC-generated diagrams displayed together for readability. Graphviz renders
the diagrams; terminal color-control sequences are omitted.

\begin{center}
\small
\begin{tabular}{lrrrl}
\toprule
Representation & Logic nodes & Edges & Levels & Function statistic\\
\midrule
Original SOP & 8 & 16 & 3 & \texttt{cube = 10}\\
Local AIG (\texttt{aig}) & 8 & 16 & 3 & \texttt{aig = 12}\\
Strashed AIG & -- & -- & 4 & \texttt{and = 12}\\
Local BDD (\texttt{bdd}) & 8 & 16 & 3 & \texttt{bdd = 16}\\
Collapsed BDD & 4 & 14 & 1 & \texttt{bdd = 14}\\
SOP after \texttt{strash; logic} & 12 & 24 & 4 & \texttt{cube = 12}\\
\bottomrule
\end{tabular}
\end{center}
The function counters are ABC's reported values, not interchangeable measures
of network size. Every case has four inputs, four outputs, and no latches.

\newpage
\subsection*{Reading and displaying the original network}
\begin{verbatim}
read mul.blif
print_stats
show
\end{verbatim}
The original network has eight internal nodes: four partial-product gates,
two XOR gates, and two carry gates. XOR nodes have two cubes, so the total
is ten cubes. Its longest path has three logic levels.
\evidence{ex2-read}{0.73}{Original SOP network and its measured statistics.}

\newpage
\subsection*{Structural hashing}
\begin{verbatim}
read mul.blif
strash
print_stats
show
\end{verbatim}
\texttt{strash} constructs a global AIG with two-input AND nodes and
complemented edges. Each XOR requires a small AND/inverter implementation.
The resulting network has twelve AND nodes and four levels. More AIG
nodes than original logic nodes is not a correctness problem: the node
types represent different amounts of logic.
\evidence{ex2-strash}{0.73}{Structurally hashed AIG. Inversion is represented on edges.}

\newpage
\subsection*{Collapsing and displaying all output BDDs}
\begin{verbatim}
read mul.blif
strash
collapse
print_stats
show_bdd -g
\end{verbatim}
After collapse, there are four logic nodes, one per output, and one logic
level. Their functions are expressed directly in terms of primary inputs.
The network has fourteen fanin edges, and ABC reports
\texttt{bdd = 14}. The \texttt{-g} option displays all output functions;
without it, \texttt{show\_bdd} defaults to one output's driver.
\evidence{ex2-collapse}{0.72}{Global view of all four multiplier output functions.}

\newpage
\section{Exercise 3: representation comparisons}
\subsection{Local AIG versus structurally hashed AIG}
\begin{verbatim}
read mul.blif
aig
print_stats
show
\end{verbatim}
\texttt{aig} preserves the eight-node logic network and replaces each
node's local function representation by an AIG. Its fanins are still the
signals feeding that logic node. In contrast, \texttt{strash} constructs
one global AND/inverter network whose leaves are the circuit inputs.

For this multiplier, \texttt{aig} reports twelve ANDs in the local function
representation, while \texttt{strash} also produces twelve AND nodes.
There is no observed reduction in that AND count. The structural difference
is visible in the topology: eight logic nodes and three levels versus
twelve global AND nodes and four levels. Structural hashing merges
structurally identical nodes; it is not arbitrary Boolean minimization.
Local AIGs should not be described simply as ``unhashed AIGs.''
\evidence{ex3-aig}{0.65}{Local AIG representation; compare with the global AIG on the previous pages.}

\newpage
\subsection{Local BDD versus collapsed BDD: network structure}
\begin{verbatim}
read mul.blif; bdd; print_stats; show
read mul.blif; collapse; print_stats; show
\end{verbatim}
\texttt{bdd} changes local node functions without flattening the network:
eight nodes, sixteen edges, and three levels remain. Its
\texttt{bdd} counter is sixteen. \texttt{collapse} composes the functions
into output functions over primary inputs, leaving four nodes, fourteen
edges, and one level, with a reported BDD counter of fourteen.

These observations distinguish the circuit's structure from the size of its
function representations. Collapse reduces these measurements for this
example; it need not make BDDs small for an arbitrary circuit or ordering.
The diagram regions below are cropped from the full screenshots to keep
the node labels readable.
\comparison{ex3-bdd-diagram}{Local BDD network: internal signals remain as fanins.}
{ex3-collapse-diagram}{Collapsed network: each output function depends directly on primary inputs.}

\newpage
\subsection*{Local and global functions for \(y_2\)}
\begin{verbatim}
read mul.blif; bdd; show_bdd y2
read mul.blif; collapse; show_bdd y2
\end{verbatim}
Before collapse, \(y_2=p_{11}\oplus c_1\) is a local two-variable
function. After collapse, the intermediate signals are substituted by
their definitions, so the BDD tests primary inputs instead.
Using \texttt{show\_bdd -g} for both cases would instead request a global
view in both cases and would not expose this local-versus-global distinction.
\comparison{ex3-local-bdd-diagram}{Local BDD for \(y_2\), in terms of its two immediate fanins.}
{ex3-global-bdd-diagram}{Collapsed BDD for \(y_2\), in terms of primary inputs.}

\newpage
\subsection{Converting a strashed AIG to SOP logic}
\begin{verbatim}
read mul.blif
strash
logic
print_stats
lsv_print_nodes
\end{verbatim}
\texttt{logic} converts the strashed AIG into a logic network with SOP
node functions. It is the necessary representation conversion; applying
\texttt{sop} directly to a strashed network is not sufficient because
\texttt{sop} requires a logic network.

The resulting network has twelve nodes, twenty-four edges, twelve cubes,
and four levels. The \texttt{lsv\_print\_nodes} listing confirms SOP data,
including \texttt{11 1} for AND nodes and complemented-input cubes such
as \texttt{01 1}. Some functions use an OFF-set cube such as
\texttt{00 0}; a zero-valued cube is valid and should not be mistaken
for a missing function. Conversion preserves the represented function,
but does not reconstruct the original eight-node XOR-level design.
\evidence{ex3-sop}{0.60}{Beginning of the actual SOP node listing after conversion.}
\end{document}
